\SourceLeaf{320}%
zugrunde, die in der vorliegenden Untersuchung erfüllt ist. Man muss nur beachten, dass \(\psi^0\) und \(\psi'\) dort genau \(T+\frac12t-(U+\frac12u)\) und \(T-\frac12t-(U-\frac12u)\) hier entsprechen, dass \(h\) dort dasselbe wie \(s\) hier bedeutet und schließlich, dass sich das dortige \(\chi\) von der hier mit \(U\) bezeichneten Größe im Allgemeinen nur um eine Größe zweiter Ordnung unterscheiden kann, der Unterschied aber in jedem Fall mindestens von erster Ordnung ist. So ergibt sich bis zur dritten Ordnung einschließlich genau
\[
\begin{aligned}
T+\tfrac12t&=U+\tfrac12u-\frac{ee\cos P\mathord{.}\sin P\mathord{.}\sin U\mathord{.}\cos U^2}{12\cos\varphi\cos\theta}\mathord{.}s^3\\
T-\tfrac12t&=U-\tfrac12u+\frac{ee\cos P\mathord{.}\sin P\mathord{.}\sin U\mathord{.}\cos U^2}{12\cos\varphi\cos\theta}\mathord{.}s^3
\end{aligned}
\]
und folglich mit derselben Genauigkeit
\[
\begin{aligned}
T&=U\quad .\ .\ .\ .\ .\ .\ .\ .\ .\ .\ .\ .\ .\ .\ .\ .\ .\ .\ .\ .\quad(6)\\
t&=u-\frac{ee\cos P\mathord{.}\sin P\mathord{.}\sin U\mathord{.}\cos U^2}{6\cos\varphi\cos\theta}\mathord{.}s^3\quad .\ .\ .\ .\ .\ .\ .\ .\ .\quad(7)
\end{aligned}
\]

\indent Der Vergleich der Länge der geodätischen Linie auf dem Ellipsoid mit dem Großkreisbogen auf der Kugel wurde in Artikel 15 für den hier betrachteten Fall zwar nicht eigens ausgeführt, lässt sich aber sehr leicht ergänzen. In der dortigen Schreibweise ist die Länge des geodätischen Bogens
\[
=A\int\frac{\cos y}{m\cos\psi}\mathord{.}\mathrm{d}x
\]
wobei die Integration von \(x=-\frac12(h-\delta)\) bis \(x=+\frac12(h+\delta)\) reicht. Da \(y\) und \(\psi\) nur Größen dritter Ordnung sind, sieht man leicht, dass das Weglassen des Faktors \(\frac{\cos y}{\cos\psi}\) im Wert des Integrals nur einen Fehler siebter Ordnung hervorrufen kann. Die Länge ist also bis zur fünften Ordnung einschließlich genau
\[
\begin{aligned}
=A\int\frac{\mathrm{d}x}{m}&=A\int\mathrm{d}x(1-\mu x^3-\mu'x^4)\\
&=A(x-\tfrac14\mu x^4-\tfrac15\mu'x^5)+\text{Const.}\\
&=A\{h-\tfrac18(h^3\delta+h\delta^3)\mu-\tfrac1{80}(h^5+10h^3\delta\delta+5h\delta^4)\mu'\}
\end{aligned}
\]
Die Koeffizienten \(\mu,\mu'\) lassen sich angeben, indem man in der Reihe
\[
m=1-\frac{2ee\cos P\mathord{.}\sin P}{3\cos\varphi\cos\theta}\mathord{.}q^3-\frac{ee\cos P^2}{6\cos\varphi^2\cos\theta^2}(1-7ee\sin P^2)q^4\quad .\ .\ .\ .
\]
(die sich unmittelbar aus der in Artikel 9 angegebenen Reihe ergibt) für \(q\) die Substitution
